% Nouveau EGA V (cote 25): archive cote 25, site page 70 (PDF page 71).
% typescript. Draft compared: none.
% AI-transcribed by Claude Opus 5.5 (claude-opus-5-5, adaptive thinking, effort high), Message Batches API, batch ega25b (msgbatch_01A7n5YuXubUPNK1zsepi8RG), 2026-11-13; not yet checked by eye.
\scanpage{25/70}
\mspage{65}

est bien indépendant de la base\hand{/}choisie.  Dans le cas  m=1, on a
l'interprétation intrinsèque évidente signalée dans 5.2.  Avec cette
convention de langage, 5.3.  reste valable t el quel, d e même 5.4.
La première partie de la remarque 5.5.   admet une \xout{généralisation} \inserted[typed]{\xout{variante pour le}}
\inserted[typed]{généralisation au}
\xout{au} cas m quelconque:  lorsque $F_s$  est \xout{\unsure{Sm}} \inserted[typed]{$(S_m)$}, alors la condition de
régularité envisagée pour $\phi_{\hand{$\xi$}}$  s'exprime de façon purement dimensionelle.
La deuxième partie de la remarque 5.5. vaut telle quelle dans le cas
m quelconque.      Théorè\unsure{m|p}e 5.6.  s'étend tel quel,  de même 5.7.

\uline{Proposition} 6.1.  :  lire   dim $f(X_i) \hand{$\geqslant$}$  m+1 ,  et
plus loin   dim $f(Z) \hand{$\leqslant$}$ m-1.

Les laius généraux d e 7.1. sont valables tels quels dans le cas
m quelconque.  7.2. et 7.3.  mutatis mutandis (attention dans 7.2.
à la not\unsure{a}tion m, confusante ici),  d'autre part dans la démonstration
de 7.4. on n'a plus besoin de procéder de proche en proche, mais on
peut prendre directement une section linéaire  de codimension m=n .

Dans 8.2. , remplacer l\unsure{a}\xout{hypothèse} condition    dim $Y_{\hand{$\xi$}}$ \hand{$>$} dim X-\unsure{1}
par  dim $Y_{\hand{$\xi$}}$ \hand{$>$}  dim X  - \over{1}{m} , et l'hypothèse  dim $f(X_i) \hand{$>$} 0$  par
dim $f(X_i) \hand{$\geqslant$}$ m.   Modification analogue dans la suite de 8.2.  Comme
8.3.  présente un exemple, il n'y a pas lieu de le changer, donc on
y gardera m=1.   Je te laisse le soin de faire l'exercice,demandant
sans doute\hand{/}un peu plus de soin, de trouver  les bons énoncés pour m
quelconque correspondants à  8.4.  8.5.  et 8.6.   (pages \inserted{30,}31,32 )\hand{;}
\hand{il n'est d'ailleurs sans doute pas indispensable de}
\hand{faire cet exercice, s'il te barbe de le faire.}

\mspage{70}
